Este trabajo tiene como alcance la exploración de modelos y distancias en un mismo sistema de aprendizaje automático orientado a la extracción de los elementos sonoros de una pieza musical: voz, bajo, batería y acompañamiento, junto con la implementación de una interfaz gráfica que permite hacer diferentes pruebas. 

El trabajo abarca la implementación de distintas funciones de pérdida (\textit{L1, L2}, pérdidas en frecuencia, pérdidas de máscara\dots), el entrenamiento controlado de modelos con una misma configuración de \textit{hiperparámetros}, y la evaluación del rendimiento obtenido a través de métricas específicas.

El proyecto no pretende competir con modelos de última generación en separación musical, sino explorar el impacto de diferentes estrategias de entrenamiento y funciones de pérdida en arquitecturas más simples, trabajando con \textit{datasets} ya creados y recursos de cómputo reducidos.